% Generated draft: requires mathematical and source review.
\documentclass[11pt]{article}
\usepackage[margin=25mm]{geometry}
\usepackage{fontspec}
\setmainfont{Latin Modern Roman}
\usepackage{amsmath,amssymb,mathtools,unicode-math}
\setmathfont{Latin Modern Math}
\usepackage{enumitem,needspace,xurl,hyperref}
\setlength{\parindent}{0pt}
\setlength{\parskip}{5pt}
\setlength{\emergencystretch}{3em}
\newcommand{\N}{\mathbb{N}}
\newcommand{\Z}{\mathbb{Z}}
\newcommand{\Q}{\mathbb{Q}}
\newcommand{\R}{\mathbb{R}}
\newcommand{\C}{\mathbb{C}}
\newcommand{\F}{\mathbb{F}}
\newcommand{\A}{\mathbb{A}}
\newcommand{\PP}{\mathbb{P}}
\newcommand{\E}{\mathbb{E}}
\newcommand{\eps}{\varepsilon}
\newcommand{\dd}{\,\mathrm{d}}
\begin{document}
\section*{UnsolvedMath ID: 30003107}
\textbf{Source problem code:} catalogue entry
\subsection{Definitions and mathematical statement}
Let \(K\subset\overline{\Q}\) be a number field, let \(X\) be a smooth projective variety over \(K\), and let \(f\) be a dominant self-map defined over \(K\). The supplied record does not specify whether a self-map must be a morphism \(f\colon X\to X\) or may be a rational map \(f\colon X\dashrightarrow X\). Both interpretations are described below; the intended scope requires confirmation.

For the rational-map interpretation, write \(I_f\) for the indeterminacy locus and require that \(P\in X(\overline{\Q})\) have a well-defined forward orbit: starting with \(P_0=P\), each \(P_n\) lies outside \(I_f\), and \(P_{n+1}=f(P_n)\) for every integer \(n\geq 0\). Write
\[
\mathcal O_f(P)=\{P_n:n\geq 0\}.
\]
For a morphism, every algebraic point has a well-defined forward orbit. Zariski density means
\[
\overline{\mathcal O_f(P)}^{\mathrm{Zar}}=X_{\overline{\Q}}.
\]

Choose an ample Cartier divisor \(H\) on \(X\) and an associated logarithmic Weil height \(h_H\colon X(\overline{\Q})\to\R\). Set \(h_H^+(Q)=\max\{1,h_H(Q)\}\). The arithmetic degree, when its defining limit exists, is
\[
\alpha(f,P)=\lim_{n\to\infty}\bigl(h_H^+(P_n)\bigr)^{1/n}.
\]
For \(d=\dim X\geq 1\), interpret the dynamical degree in the record as the first dynamical degree
\[
\delta(f)=\lim_{n\to\infty}\left(\bigl((f^n)^*H\bigr)\cdot H^{d-1}\right)^{1/n}.
\]
Here the intersection number is taken on \(X\); for rational maps, pullback is defined through the graph of the indicated iterate \(f^n\). In particular, the formula uses \((f^n)^*\), without assuming it equals \((f^*)^n\). These are the usual height and intersection definitions; the resulting degrees, when defined as above, do not depend on the ample divisor or associated height chosen.

\textbf{Conjectural statement.} For every such \(K,X,f\) and every algebraic point \(P\) with a well-defined Zariski-dense forward orbit, does the arithmetic-degree limit exist and satisfy
\[
\alpha(f,P)=\delta(f)?
\]
Under the morphism interpretation, the well-defined-orbit condition is automatic. Under the rational-map interpretation, it is a necessary explicit qualification absent from the supplied wording. The record also does not state whether existence of the arithmetic-degree limit is assumed or is part of the conjectural conclusion; the displayed formulation includes existence in the conclusion. The intersection formula above treats positive-dimensional varieties; the source does not specify a convention for dimension zero.

The supplied attribution identifies this as the Kawaguchi--Silverman conjecture. Its supplied dated assessment labels the problem partially solved, says that the abelian-variety case is proved, mentions further polarized, int-amplified, and other cases, and says that the general dominant-map claim remains open. Those are supplied curation claims, not mathematical or literature findings independently verified here; the hypotheses of the cited special cases are not supplied.
\subsection{Short English statement}
The conjecture asks whether every algebraic point with a Zariski-dense forward orbit has height growth rate equal to the first dynamical degree of the map. The source leaves unresolved whether rational maps are included and whether existence of the height-growth limit is assumed. Rational maps require every iterate of the point to be defined. The supplied partially solved label and descriptions of known cases require independent review.
\subsection{Sources}
\begin{itemize}
\item Supplied local UnsolvedMath registry record, permanent ID \(30003107\), title \emph{Kawaguchi--Silverman Arithmetic Degree Conjecture}, Oberwolfach Reports Open Problems collection; exact source problem code \texttt{catalogue entry}. The record gives proposed year 2016 and names no individual proposer.
\item Supplied report citation: \emph{Diophantische Approximationen (2017)}, Oberwolfach Reports, DOI \texttt{10.4171/owr/2016/21}, pp.\ 1099--1187. \url{https://doi.org/10.4171/owr/2016/21}. The background describes Oberwolfach Reports as documenting workshops and their open-problem discussions. The supplied title year and the year appearing in the DOI have been retained without reconciliation.
\item Supplied literature source, labeled ``Kawaguchi--Silverman'': \url{https://arxiv.org/abs/1208.0815}.
\item Supplied literature source, labeled ``Silverman on abelian varieties'': \url{https://www.numdam.org/item/JTNB_2017__29_1_151_0/}.
\item Supplied literature source, labeled ``Meng--Zhong'': \url{https://arxiv.org/abs/2408.00566}.
\item Literature assessment date: \texttt{2026-08-21}. This is supplied dated curation, described in the record as literature triage derived from the supplied OWR curation, and was not independently verified here. It labels the problem partially solved, reports the abelian-variety case as proved, mentions polarized, int-amplified, and other cases, and reports the general dominant-map claim as open. The record expressly calls for independent verification of time-sensitive status and cited claims.
\end{itemize}
\end{document}
